\begin{center}
    \section*{Abstract}
\end{center}

% Use one-and-a-half spacing locally so the abstract fits on one page.
% Remove paragraph indentation; use a small vertical skip between paragraphs instead.
\begin{onehalfspace}
\setlength{\parindent}{0pt}
\setlength{\parskip}{4pt}

\large
Getting timely, trustworthy medical care is harder than it should be. People ignore symptoms until things get serious, or they search online and come away more confused. Formal access --- finding a doctor, booking, paying, actually consulting --- involves so many disconnected steps that many just give up. Doctors, meanwhile, spend a disproportionate chunk of their time on paperwork rather than patients. EHealth is our attempt to close that gap. It is an AI-assisted telemedicine platform that ties together four things that normally exist in isolation: intelligent symptom triage, doctor discovery and booking with integrated payment, live WebRTC video consultation, and blockchain-verified prescription records.  The triage module accepts free-form text or voice input --- grammatically broken or vague input included --- and returns a structured JSON response classifying urgency as HIGH, MEDIUM, or LOW, plus first-aid guidance and a hospital-lookup flag. Powered by the Baichuan-M2-32B LLM with constrained prompt engineering, it achieved \textbf{87.5\%} overall accuracy, \textbf{94.8\%} sensitivity on critical cases, and zero critical false negatives on a 200-prompt test set. After a consultation ends, the platform automatically picks up the recorded audio, transcodes it to 16\,kHz mono WAV via FFmpeg, transcribes it locally with \textit{Xenova/whisper-tiny}, and passes the transcript to the LLM to produce a structured clinical draft (symptoms, diagnosis, medications, follow-up). This is rendered as a PDF and pushed to Google Drive. On 50 test dialogues, the pipeline reached a \textbf{14.2\%} general WER (8.6\% on medical terms) and a mean \textbf{F1-score of 93.0\%} across all four clinical fields. The prescription verification module anchors the SHA-256 hash of every finalized prescription to a smart contract on the Polygon Proof-of-Stake network. Patients can re-upload a downloaded PDF and get a dual answer --- database match and on-chain match --- rather than a single boolean. Across nine test transactions over six weeks, the median on-chain write cost was \textbf{0.00607\,POL} (${\approx}$\$0.0005), proving the integrity layer adds negligible overhead at any scale.

\vspace{3mm}
\textit{\textbf{Keywords:} Telemedicine, Medical Triage, Large Language Models, Automatic Speech Recognition, Blockchain, Smart Contracts, Polygon PoS, Prescription Verification, NestJS, Flutter, Clinical Decision Support.}

\end{onehalfspace}